\subsection{Direct Sums and Complements}

Throughout this subsection, $R$ denotes a nonzero commutative ring with identity and $M$ denotes an $R$-module.

\begin{definition} \label{def:module-direct-sums} \hlt{Direct Sums and Direct Summands}
\begin{enumerate}[label=(\roman*)]
\setlength{\itemsep}{0pt}
\item For $R$-modules $M_1,\ldots,M_n$, their \hlt{external direct sum} is $\bigoplus\limits_{i=1}^{n}M_i\defeq\{(m_i)_{i=1}^{n}:m_i\in M_i\}$\\
with componentwise $R$-module operations.
\item For submodules $\{S_i:i\in I\}$ of $M$, their \hlt{internal direct sum} is denoted $M=\bigoplus_{i\in I}S_i$,\\
which means $M=\sum\limits_{i\in I}S_i$ and $\forall i\in I$, $S_i\cap\sum_{j\neq i}S_j=\{0\}$.
\item A submodule $S\leq M$ is a \hlt{direct summand} if $\exists\,T\leq M$ such that $M=S\oplus T$.\\
Such a $T$ is a \hlt{complement} of $S$ in $M$.
\end{enumerate}
\end{definition}

\begin{proposition} \label{prop:module-direct-sum-characterizations} \hlt{Direct Sum Characterizations}\\
For a family $\{S_i:i\in I\}$ of distinct submodules of $M$, the following conditions are equivalent:
\begin{enumerate}[label=(\roman*)]
\setlength{\itemsep}{0pt}
\item $\sum\limits_{i\in I}S_i$ is direct, i.e. $\forall k\in I,\ S_k\cap\sum\limits_{i\neq k}S_i=\{0\}$.
\item For every $F\subseteq_{\operatorname{fin}} I$ and $x_i\in S_i$ $(i\in F)$, $\sum\limits_{i\in F}x_i=0$ $\Longrightarrow$ $x_i=0$ $\forall i \in F$.
\item Every $0\neq x\in\sum_{i\in I}S_i$ has a unique finite nonzero decomposition,\\
$\exists!\,(F,(x_i)_{i\in F})$ where $F\subseteq_{\operatorname{fin}} I$; $0\neq x_i\in S_i$; $x=\sum\limits_{i\in F}x_i$.
\end{enumerate}
\begin{proof}
\begin{enumerate}[label=(\roman*).]
\setlength{\itemsep}{0pt}
\item[$(i)$]$\Rightarrow (ii)$. If (i) holds and $\sum\limits_{i\in F}x_i=0$, then any $x_k\neq0$ gives $x_k=-\sum\limits_{i\in F\setminus\{k\}}x_i\in S_k\cap\sum\limits_{i\neq k}S_i$, a contradiction.
\item[$(ii)$]$\Rightarrow (i)$. If $x\in S_k\cap\sum\limits_{i\neq k}S_i$, write $x=\sum\limits_{i\in F}x_i$ with $F\subseteq_{\operatorname{fin}} I\setminus\{k\}$ and $x_i\in S_i$.\\
Then $x-\sum\limits_{i\in F}x_i=0$, so (ii) gives $x=0$.
\item[$(ii)$]$\Rightarrow (iii)$. Existence follows from the definition of $\sum_i S_i$ after deleting zero terms.\\
If $\sum\limits_{i\in F}x_i=\sum\limits_{j\in G}y_j$, $0\neq x_i\in S_i$, $0\neq y_j\in S_j$, then, over $H=F\cup G$, missing terms set to $0$, $\sum\limits_{h\in H}(x_h-y_h)=0$.\\
By (ii), $x_h=y_h$ for all $h\in H$, hence $F=G$ and the terms agree.
\item[$(iii)$]$\Rightarrow (ii)$. If $\sum_{i\in F}x_i=0$, set $F_0\defeq\{i\in F:x_i\neq0\}$.\\
If $F_0\neq\emptyset$, choose $k\in F_0$. If $F_0=\{k\}$, then $x_k=0$, contradiction; otherwise $0\neq x_k=-\sum\limits_{i\in F_0\setminus\{k\}}x_i$,\\
giving two finite nonzero decompositions of $x_k$, contradiction. Hence $F_0=\emptyset$.
\end{enumerate}
\hlt{Structural Commentary}: Directness is the absence of hidden cancellation among different summands.\\
\hlt{Key Insight}: Unique decomposition is controlled by the zero vector.
\end{proof}
\end{proposition}

\begin{example} \label{ex:z-submodules-mostly-not-complemented} \hlt{Most Submodules of $\mathbb{Z}$ Are Not Complemented}\\
As a $\mathbb{Z}$-module, every submodule of $\mathbb{Z}$ is $n\mathbb{Z}$ for a unique $n\geq0$. For $m,n>0$, $m\mathbb{Z}\cap n\mathbb{Z}=\operatorname{lcm}(m,n)\mathbb{Z}\neq\{0\}$.\\
Thus $m\mathbb{Z}\oplus n\mathbb{Z}\neq\mathbb{Z}$ for all $m,n>0$, so the only complemented submodules of $\mathbb{Z}$ are $\{0\}$ and $\mathbb{Z}$.
\end{example}

\begin{theorem} \label{thm:complements-and-quotients-modules} \hlt{Complements and Quotients}\\
If $M=S\oplus T$, then $\pi|_T:T\xrightarrow{\sim}M/S$, $t\mapsto t+S$, where $\pi:M\twoheadrightarrow M/S$ is the quotient map.\\
Hence any two complements of $S$ are isomorphic.
\begin{proof} $\ker(\pi|_T)=T\cap S=\{0\}$. Also, for $m=s+t$ with $s\in S$, $t\in T$, one has $m+S=t+S$; hence $\operatorname{im}(\pi|_T)=M/S$. Thus $\pi|_T$ is an isomorphism, and all complements are isomorphic to $M/S$.\\
\hlt{Structural Commentary}: A complement gives a concrete copy of the quotient inside the original module.\\
\hlt{Key Insight}: Quotienting by $S$ remembers exactly the complementary coordinate.
\end{proof}
\end{theorem}
